\subsection{Описание предметной области}
\label{sec:analysis:domain}

Значительная часть сделок с подержанными вещами в Беларуси проходит через площадки частных объявлений. Продавец публикует описание товара и цену, покупатель находит объявление через поиск и связывается с продавцом напрямую. Сама площадка в сделке не участвует, а хранит объявления, предоставляет поиск по ним и средства связи между сторонами. Крупнейшая площадка такого рода в стране~-- \textit{Kufar.by}. На ней публикуются объявления о продаже транспорта, электроники, мебели, одежды, спортивного инвентаря, а также о продаже и аренде недвижимости [1].

Объём предложения велик даже в пределах отдельной категории. По выборке, выполненной через поисковый интерфейс площадки в сентябре 2026 года, в категории «Велоспорт» одновременно опубликовано 22\,966 объявлений. Из них 16\,494 относятся к самим велосипедам, остальные к запчастям и аксессуарам: камерам, втулкам, переключателям. В категории «Квартиры» выставлено 23\,930 объявлений о продаже и 4\,442 объявления о долгосрочной аренде. Из 13\,028 квартир, предложенных в Минске, частные лица разместили около 800, остальные 94\,\% опубликованы агентствами недвижимости. Выдача при этом не стоит на месте. Проданные товары снимаются с публикации, появляются новые объявления, продавцы меняют цены в уже опубликованных.

Со стороны спроса на площадке работают две группы пользователей. Первую составляют частные покупатели, которым нужен конкретный товар для себя: велосипед с определённым размером рамы, ноутбук с заданным объёмом оперативной памяти, квартира в выбранном районе. Покупатель ищет эпизодически, в течение нескольких дней или недель, пока не найдёт подходящее предложение. Ко второй группе относятся перекупщики, которые приобретают товар ниже рыночной цены и затем продают его дороже. Перекупщик следит за площадкой постоянно и сразу по нескольким категориям, а доход перекупщика прямо зависит от скорости обнаружения выгодного объявления. Интенсивность поиска у двух групп разная, однако потребность одна и та же: узнать о подходящем объявлении раньше других.

Рынок подержанных товаров устроен иначе, чем розничная торговля. В магазине один и тот же товар продаётся под одним артикулом, а цена меняется постепенно и примерно одинаково у разных продавцов. На площадке объявлений каждое предложение уникально. Два велосипеда одной модели различаются состоянием, комплектацией и ценой, причём цену назначает сам продавец. При спешной продаже или плохом знании рынка цена оказывается заметно ниже обычной. Такое объявление быстро находит покупателя, и выгодное предложение нередко живёт на площадке лишь несколько часов. Пользователь, который проверяет выдачу утром и вечером, регулярно пропускает объявления, выставленные и проданные в промежутке между проверками.

Ручной мониторинг площадки упирается в следующие проблемы:
\begin{itemize}
	\item объявления появляются круглосуточно, в том числе ночью, и человек не может постоянно следить за выдачей;
	\item человек не удерживает в памяти рыночный минимум по каждому интересующему товару и не всегда может сразу оценить, выгодна ли увиденная цена;
	\item при повторном просмотре большую часть выдачи, до 90\,\%, составляют уже виденные объявления, и среди них приходится заново искать несколько новых;
	\item в карточке объявления отображается только текущая цена, поэтому снижение цены на просмотренное ранее предложение легко пропустить;
	\item сравнение объявлений сразу по нескольким параметрам (цене, состоянию, характеристикам товара, району) утомительно и сопровождается ошибками.
\end{itemize}

Частично эти проблемы закрывает встроенная функция площадки: сохранённые поиски. Пользователь настраивает фильтры выдачи, сохраняет поиск в личном кабинете и получает оповещения о новых объявлениях, которые под него попадают. Возможности механизма ограничены. Цена в сохранённом поиске задаётся лишь как фильтр выдачи, снижение цены на объявление, опубликованное раньше, отдельным событием не считается, история цен не ведётся. Оповещения приходят только через сервисы самой площадки. Сторонняя программа прочитать сохранённые поиски не может: программный интерфейс площадки отвечает на такой запрос ошибкой 401, если в нём нет авторизованной сессии владельца. Подробнее этот механизм разобран в подразделе~\ref{sec:analysis:analogues}.

Если отвлечься от пользовательского интерфейса, задачу мониторинга можно описать набором объектов и запросом к ним. В предметной области выделяются следующие объекты:
\begin{itemize}
	\item пользователь~-- человек, который взаимодействует с системой через \textit{Telegram} и определяет, что именно нужно отслеживать;
	\item отслеживаемый товар~-- набор критериев одного пользователя: категория, название, фильтры по характеристикам, пороговая цена и слова-исключения;
	\item объявление~-- сведения о предложении, полученные с площадки: заголовок, описание, цена, ссылка, местоположение и дата публикации;
	\item совпадение~-- факт того, что конкретное объявление удовлетворило критериям конкретного отслеживаемого товара, вместе с типом события (новое объявление или снижение цены);
	\item уведомление~-- факт отправки пользователю сообщения о совпадении с ценой объявления на момент отправки.
\end{itemize}

К этим объектам добавляется справочник характеристик категорий площадки. Площадка не публикует перечень своих фильтров, и для построения меню выбора бренда, размера рамы или числа комнат система собирает этот справочник сама по действующим объявлениям.

Каждый цикл опроса сводится к одному и тому же запросу. Для каждого активного отслеживаемого товара нужно найти объявления, которые одновременно:
\begin{itemize}
	\item имеют цену строго ниже пороговой;
	\item соответствуют категории и фильтрам товара и не содержат слов-исключений;
	\item опубликованы после того, как товар был добавлен в систему, либо подешевели с момента предыдущего опроса;
	\item ещё не отправлялись этому пользователю по такой же или более низкой цене.
\end{itemize}

Первые два условия проверяются по свежему ответу площадки. Третье и четвёртое сравнивают этот ответ с тем, что система сохранила раньше: без хранимой истории объявлений и отправленных уведомлений их проверить нельзя. Объявления, прошедшие отбор, фиксируются как совпадения и ставятся в очередь на доставку. При интервале опроса в 3 минуты система выполняет 480 циклов в сутки. В каждом цикле для каждого товара приходит до 42 объявлений, и каждое из них сверяется с базой данных. Устройство системы определяет этот поток сравнений, а отображение сообщений в нём вторично.
